% !TeX root = ../main.tex
\section{Comparison, failure modes and open problems}
\label{sec:comparison}

\subsection{Comparison}

The commercial case needs three properties at once: no dry signal of the target,
an output a real-time chain can use, and robustness to separated stems.
Table~\ref{tab:methods} shows that no published method has all three. Effect
recognition returns usable parameters but was trained and tested almost only on
synthetic, isolated renders. Unpaired modelling needs no aligned dry signal and
returns a real-time model, but covered only distortion on isolated guitar. Tone
transfer (EG-VAE) handles whole chains but returns audio. Inference-time search
(ST-ITO) drives any chain but fails on guitar tone with an embedding objective.
StemFX is the only work trained on separated stems, for mixing effects.

\begin{table}[h]
\centering
\caption{Main methods by requirement and output. ``Non-diff.'': works with a
chain that cannot be differentiated. ``Sep.'': trained or tested on separated
stems. Licence of released code or weights (NC = non-commercial; -- = not
released or no licence file; ? = code released, licence not checked).}
\label{tab:methods}
\footnotesize
\setlength{\tabcolsep}{2.5pt}
\resizebox{\textwidth}{!}{%
\begin{tabular}{@{}llllllll@{}}
\toprule
Method & Needs & Output & Non-diff. & Guitar & Sep. & Code & Reported \\
\midrule
FxNet~\cite{comunita2021guitar} & wet & drive type, gain, tone & yes & yes & no & BSD-3 & 81--91\%, MAE $\approx$0.04 \\
Hinrichs 2022~\cite{hinrichs2022convolutional} & wet in mix & type, params & yes & yes & no (mix) & -- & 90.0\% in mixes \\
Okita~\cite{okita2026audio} & wet & chain (search) & yes & yes & no & -- & 23.07\dB{} wet SI-SDR \\
StemFX~\cite{cheng2026stemfx} & wet stems & mixing FX chain & yes & no & yes & MIT & 86.8\% top-1 \\
DeepAFx-ST~\cite{steinmetz2022style} & wet + input & EQ, comp. & SPSA & no & no & NC & -- \\
ST-ITO~\cite{steinmetz2024stito} & wet + input & any chain (search) & yes & fails & no & Apache-2.0 & 0.86 style acc. \\
TAPE~\cite{long2026tape} & wet + input & own DSP chain & no & partly & no & -- & STFT L1 1.12 vs 2.34 \\
Fx-Encoder++~\cite{yeh2025fxencoder} & mix + query & embedding & -- & partly & yes & NC & R@10 39.4 $\to$ 20.5\% \\
Unpaired GAN~\cite{wright2023adversarial} & wet + other DI & neural amp & -- & yes & no & -- & MUSHRA 77--84 \\
Diffusion blind~\cite{moliner2025unsupervised} & wet + clean prior & neural / grey-box & -- & yes & no & ? & stable from 18\,s \\
Remove + model~\cite{hinrichs2024blind} & wet & neural effect & -- & yes & no & -- & MUSHRA $\geq$92 (1 effect) \\
EG-VAE~\cite{yeh2026egvae} & wet + any DI & audio, DI estimate & -- & yes & no & -- & MOS 4.30 / 3.70 \\
Tone embedding~\cite{chen2024zeroshot} & wet ref.\ + DI & neural amp & -- & yes & no & -- & 2$\times$ loss unseen \\
Open-Amp~\cite{wright2025openamp} & paired (enrol) & neural, 394 devices & -- & yes & no & -- & 91.1\% class acc. \\
Tone Match~\cite{fractalndtonematch} & isolated wet & EQ / IR & yes & yes & no & product & -- \\
\midrule
This app & stem + notes & app chain + IR & yes & yes & \wip & GPL-3.0 & Section~\ref{sec:results} \\
\bottomrule
\end{tabular}}
\end{table}

\subsection{Known and expected failure modes}

\begin{description}[style=nextline,leftmargin=1.2cm]
\item[Separation artefacts] Effect retrieval roughly halves from isolated to
  separated stems \reported~\cite{yeh2025fxencoder}. Reverb tails, stereo delays
  and wide choruses are the most exposed to leakage and masking; bleed from keys
  or other guitars adds energy a matcher will try to reproduce \expected.
\item[Several guitars in one stem] Double-tracked or layered guitars give one
  stem with mixed tones. The pan split separates parts on different sides, not
  in the same position \expected.
\item[Mixing and mastering] Bus compression, limiting and mix EQ act after the
  guitarist's chain. An estimator attributes them to the guitar; a final linear
  match absorbs the EQ but not the dynamics \expected.
\item[Non-uniqueness] Drive vs amp gain, fuzz vs high-gain amp, compression vs
  picking dynamics, effect order. Device identity is confused even between clean
  renders of similar circuits~\cite{comunita2021guitar}, and order is harder than
  presence~\cite{wada2025hyperbolic}. Harmless for reproduction if the result
  sounds alike, which is why audio objectives are preferred~\cite{peladeau2024blind}.
  The app's benchmark shows it: the right amp voicing is found in only 7 of 12
  items, yet the sound gets closer \measured.
\item[Equipment outside the chain] Guitar, pickups and microphone shape the tone
  and can only be compensated by EQ. A mono chain cannot reproduce stereo
  effects. Generic blocks may lack a specific amp's character, as EG-VAE's
  differentiable baseline showed~\cite{yeh2026egvae}.
\item[Search pathologies] LFO rates and delay times make rugged
  landscapes~\cite{okita2026audio,uzrad2024diffmoog}; an embedding objective can
  improve while the sound moves away~\cite{long2026tape}; errors of a removal
  model add to errors of the estimate~\cite{take2024audio}.
\item[Tone changes within a song] Clean verses and distorted choruses need
  different presets. Segmenting songs into tone sections exists only in
  products~\cite{bouharb2026gptonebuilder}.
\end{description}

\subsection{Open problems}

\begin{enumerate}
\item Estimating an interpretable amp, cabinet and effects chain from a separated
  stem of a commercial song, evaluated against known equipment.
\item A benchmark for that task: commercial-like mixes with known chains, and
  real multitracks with the guitarist's DI and amp tracks, scored with metrics
  validated by listening tests.
\item Unpaired or blind modelling of a full chain, on separated rather than
  isolated guitar.
\item Delay, reverb and modulation estimation on real guitar phrases, with
  validated metrics.
\item Training data that can be redistributed: most realistic sets and the
  largest capture libraries are non-commercial or forbid bulk use.
\end{enumerate}
